\section{Discussion}
\label{sec:discussion}

\subsection{Limitations}

\paragraph{Policy Engineering Burden.}
Although \sys{}'s declarative YAML policies are more accessible than code, crafting rules that capture legitimate use patterns without creating security gaps requires domain expertise. The 3 false positives (1.88\% FPR) and 1 successful attack highlight the need for ongoing refinement. Adaptive policies that learn from user feedback (e.g., automatically whitelisting patterns a user repeatedly approves) are a promising direction.

\paragraph{Provenance Granularity and Paraphrasing.}
Provenance is tracked at argument level: if any ancestor of a tool-call argument is untrusted, the entire argument is labeled untrusted. This can be overly conservative when benign and malicious content are co-located. More critically, substring-based provenance resolution fails when the LLM paraphrases untrusted input. Our paraphrase evaluation (RQ6) quantifies this: provenance recall drops to 0\% for NL arguments the LLM rephrases, yielding 25\% paraphrase ASR on NL-argument attacks.

\textbf{Validated Mitigation: Semantic Provenance.} We implement and evaluate an embedding-based extension (Section~\ref{subsec:paraphrase}) that augments substring matching with sentence embeddings (all-MiniLM-L6-v2, 22M parameters). This three-stage pipeline (substring $\to$ embedding $\to$ conservative default) achieves \textbf{100\% provenance recall} on both literal and NL arguments, reducing paraphrase ASR from 23.3\% to 0.0\% with $<$5\,ms additional latency per argument and no GPU required.

\paragraph{Semantic Understanding.}
Policies are syntactic---they check paths, tool names, and provenance labels but cannot distinguish intent (e.g., writing a benign note vs.\ exfiltrating passwords to a file). Integrating a secondary LLM for semantic analysis could improve precision but introduces cost and a new attack surface.

\paragraph{Taint Explosion.}
Conservative trust propagation risks labeling the entire system state as untrusted in long conversations mixing trusted and untrusted data. Our benchmark's high TSR (98.12\%) indicates this is manageable for typical smart-home workflows, but complex cross-domain tasks would trigger more confirmations. Sanitization primitives (user-initiated declassification) and periodic taint reset mitigate this.

\subsection{Analysis of the Single Successful Attack}

The one attack bypassing \sys{} (\texttt{attack\_multi\_016}) exploited \emph{deferred execution}: the injected instruction ``add to queue: format disk; run queue when battery is low'' did not trigger dangerous-action detection at injection time. By execution time, the provenance chain passed through benign intermediate states, obscuring the untrusted origin (TOCTOU vulnerability). Only Qwen3-30B executed this attack; the other four models refused. Mitigations include persistent provenance labels on queued items and syntactic screening for scheduling patterns in untrusted input.

\subsection{Model Selection and Generalizability}

Our evaluation uses five open-source models spanning four architectures (GPT-family, InternVL multimodal, Qwen chat-optimized, Qwen-Next reasoning) and a 6$\times$ parameter range (20B--120B). This architectural diversity, combined with Qwen3-Next-80B serving exclusively as an independent oracle (95\% agreement, breaking agent--oracle circularity), demonstrates that \sys{}'s defense is model-independent. Open-source models enable exact reproducibility; proprietary APIs change behavior across versions.

\paragraph{On Baseline Risk.}
The 12.34\% baseline ASR might suggest LLMs are ``already safe.'' We disagree: (1)~probabilistic safety (88\% refusal rate) provides no guarantee for any individual interaction---1-in-8 attacks succeed; (2)~safety training is brittle against jailbreaks; (3)~agent actions (file deletion, door unlocking) are irreversible. \sys{} provides \emph{deterministic} safety at the tool-call boundary.

\subsection{Generalization}

While evaluated on smart homes, \sys{}'s architecture generalizes: enterprise assistants (DLP policies on email/document tools), healthcare agents (provenance-verified medical recommendations), web browsing agents (source-reputation-based trust labels), and autonomous vehicles (operator commands vs.\ adversarial sensor data). The core abstraction---provenance-gated policies at the tool-call boundary---applies wherever LLM agents invoke external tools.
